% ---------------------------------------------------
% Trabajo Final de Grado
% Author: Eric Ríos Hamilton <alu0101549835@ull.edu.es>
% Chapter: Goals 
% ----------------------------------------------------


\chapter{Goals} \label{chap:Goals}
This project's main goal, the creation of a modular procedural terrain generation tool, can be subdivided into the following objectives, which cover the system’s main functional, architectural, and usability aspects:
\begin{enumerate}
    \item \textbf{Procedural generation of heightmaps}.
    The system must provide a variety of algorithmic approaches to heightmap generation, including noise-based methods, custom processors and compositing of multiple sources, offering designers a flexible starting point for terrain creation.
    \item \textbf{Generation of 3D geometry from a heightmap}.
    The system must be capable of converting a two-dimensional heightmap into renderable three-dimensional terrain geometry.
    This includes the generation of vertex positions, normals, and tangents required for correct lighting and shading.
    \item \textbf{Implementation of agents for terrain modification}.
    Following the agent-based approach proposed by Doran and Parberry \cite{Doran:2010:CPT}, the system must support a pipeline of software agents capable of introducing specific geographic features in a controlled and parametrizable manner, addressing the limitations of purely stochastic methods.
    \item \textbf{Real time Level of Detail (LOD) system}. 
    To support large-scale terrain within the constraints of real-time rendering, the system must implement a level of detail mechanism that reduces geometric complexity for terrain distant from the viewer while maintaining visual fidelity where it matters, without interrupting gameplay.
    \item \textbf{Modularity and fine-grained parametrization}.
    Every component of the system must be designed to be independently configurable and replaceable. 
    This ensures that the tool can adapt to the requirements of different game projects without requiring modifications to core logic.
    \item \textbf{Benchmarking system}.
    A dedicated framework must be developed to measure the computational performance of each stage of the generation pipeline.
    This allows for informed optimization decisions and provides a quantitative basis for evaluating the system's performance.
    \item \textbf{Designer-friendliness}.
    The system must be usable by game designers who are not necessarily programmers.
    Configuration of the generation pipeline should be achievable through Godot's native interface without requiring custom code, lowering the barrier to practical adoption.
    \item \textbf{Integration with Godot's native interface}.
    All generation parameters and configurations must be exposed through Godot's Resource system, enabling storage, reuse, and sharing of terrain configurations as native engine assets.

    
\end{enumerate}
